% Study Companion (BHT): Student Guide (body). Compiled by student_guide.tex.

\begin{quickstart}
\begin{enumerate}[leftmargin=1.4em]
  \item Save \texttt{BHT-Study-Companion.html} in a folder and open it with a double-click, always in the same browser.
  \item On the map of the course, click a chapter.
  \item Open an interactive example, answer \ui{Predict first}, then change one assumption at a time.
  \item In \ui{Questions}, write your justification and answer.
  \item Come back to \ui{Review} every day or two, and take a \ui{Mock exam} when a chapter feels done.
\end{enumerate}
\end{quickstart}

\section{What the app is for}

The Study Companion helps you check that you understand each chapter of the course textbook, Bhattacharya, Hyde and Tu, \emph{Health Economics} (Palgrave Macmillan, 2014), called BHT below. It covers the fifteen chapters on the syllabus. For each of them it offers models you can run and questions you can grade yourself on.

You can use it in four ways:
\begin{itemize}
  \item \textbf{Explore} the interactive examples: change an assumption and see how the equilibrium moves.
  \item \textbf{Practice} the true-or-false and multiple-choice questions of each chapter.
  \item \textbf{Review} the questions you answered, at spaced intervals, so that what you learn stays.
  \item \textbf{Test yourself} with a timed mock exam, then read your report to see which chapters are ready.
\end{itemize}
The app is a study aid. Your results stay on your computer and nobody grades them.

\subsection{The app and the textbook}

The examples and questions are written for the course; none of them is taken from BHT. They follow BHT's models and notation, and every example and every explanation tells you which section, figure or table of BHT it refers to. The app does not replace the textbook: it assumes you have read the chapter, and it does not repeat BHT's end-of-chapter exercises, which remain part of your study.

Many examples use the numbers of BHT's own worked examples. Where BHT states a model without numbers, the example uses an illustrative calibration chosen for the course. The introduction of each example says which numbers come from BHT and which are illustrative, so you know which results you should be able to find in the text.

\section{Getting started}

The app is a single file, \texttt{BHT-Study-Companion.html}. Save it in a folder on your computer and open it with a double-click.
\begin{itemize}
  \item \textbf{Browser.} It works in recent versions of Chrome, Edge, Firefox and Safari. Always open it with the same browser, because that is where your progress is kept.
  \item \textbf{No internet needed.} Everything is inside the file: text, formulas, charts and questions.
  \item \textbf{No account.} Nothing you do is sent anywhere. Your progress is saved in the browser automatically.
  \item \textbf{New versions.} When your instructor sends a new file, save it over the old one, with the same name and in the same folder, and open it with the same browser. Your progress carries over. If you are unsure, save a progress file first (Section~\ref{sg:data}).
\end{itemize}
To put your name on your report, open \ui{My report} and type it at the top.

\section{Finding your way}

The home page is a map of the course (Figure~\ref{sg:fig:map}). Chapters keep BHT's numbers and are grouped by BHT's parts. Each part has its own color, and each chapter is a tile with its number, its title and a bar that fills as you work through it.
\begin{itemize}
  \item \textbf{Colors tell you where you are.} Teal marks the opening chapter and Part~VI (Public health economics), indigo Part~I (Demand for health and health care), ochre Part~III (Information economics), plum Part~V (Health policy). Chapters of BHT that are not on the syllabus do not appear.
  \item \textbf{Green, amber and red tell you how you are doing}: right or wrong answers, and chapter readiness.
  \item \textbf{Open a chapter} by clicking its tile, or its name in the menu on the left. A chapter page has tabs: \ui{Overview}, one tab per interactive example, and \ui{Questions}. A check mark on a tab means you have explored that example.
  \item \textbf{The menu on the left} also leads to \ui{Review}, \ui{Mock exam}, \ui{My report}, \ui{Settings} and \ui{Feedback}. On a phone, use the drop-down menu at the top.
\end{itemize}
From the home page you can also continue where you left off, or open a random example you have not tried yet.

\appshot{app-map}{The map of the course. Each part of BHT has its own color; each tile is a chapter.\label{sg:fig:map}}

\section{Working with an interactive example}

Each example turns one of the chapter's models into something you can run, such as the market for lemons in Chapter~8 or the separating equilibrium of Rothschild and Stiglitz in Chapter~9. It always opens at its baseline values (Figure~\ref{sg:fig:example}).
\begin{itemize}
  \item \textbf{Assumptions} are on the left: sliders and boxes for the parameters, and buttons for alternative cases discussed in the chapter. Each slider shows its baseline value next to it.
  \item \textbf{The figure} redraws as you move a slider. The dashed gray curve shows the baseline case, so you see how far the equilibrium has moved.
  \item \textbf{The colored box} above the figure says in one line what the model concludes: whether the equilibrium holds, fails, or involves a trade-off.
  \item \textbf{The results table} shows each result now, at the baseline, and the change, with an arrow that tells you the direction.
  \item \textbf{The status line} tells you whether you are at the baseline or which assumptions you changed. \ui{Back to the baseline} resets them.
\end{itemize}
Change one assumption at a time. If you move two at once, you cannot tell which one moved the result.

Below the figure, \ui{Predict first} asks you what will happen before you touch anything. Commit to an answer: the app then sets the parameters for you and shows what the model says, with the reason. Each example has three or four of these questions. Your first answer is recorded in your report, but it does not enter readiness: a wrong prediction is part of learning a model.

\appshot{app-example}{An example after one assumption has changed: the status line names it, the dashed curve shows the baseline case, and the verdict summarizes the result.\label{sg:fig:example}}

\section{Practicing the questions}

The \ui{Questions} tab of each chapter holds 10 true-or-false statements and 20 multiple-choice questions. About half of the multiple-choice questions ask for a short calculation: keep paper and a calculator at hand.
\begin{itemize}
  \item \textbf{True or false.} Write your justification first, in two or three sentences. The \ui{True} and \ui{False} buttons unlock once you have written enough. The justification is the real exercise: a right answer with the wrong reason is still wrong.
  \item \textbf{Multiple choice.} Click the option you think is right. The wrong options are the mistakes students typically make, such as a sign error or confusing average and marginal, so a quick guess is often punished.
  \item \textbf{Your first answer counts.} You can try again as often as you like, but your report uses the first attempt. Filters at the top show the questions you have not answered or the ones you missed.
\end{itemize}

\subsection{Locked solutions}

At first the app tells you only whether your answer is right or wrong (Figure~\ref{sg:fig:questions}). The correct option and the explanation stay locked until your instructor releases them, chapter by chapter, by sending you a code. You enter each code only once, either in the box at the top of the chapter's \ui{Questions} page or in \ui{Settings}, under \ui{Solutions}. One code opens every solution of its chapter, also for the questions you already answered; a course code, if your instructor sends one, opens all chapters at once. Upper or lower case does not matter and the dash is optional.

Once a solution is visible, compare it with your justification and rate it: \ui{Yes}, \ui{Partly} or \ui{No}. Be honest; the rating is for you, and your report shows it. Each explanation cites the section of BHT to go back to.

\appshot{app-questions}{A true-or-false question answered while the solutions of the chapter are still locked.\label{sg:fig:questions}}

\section{Review, mock exam and your report}

\subsection{Review}
Every question you answer, in practice or in a mock exam, joins your review schedule. It comes back after 1~day, then 3, 7, 21 and 60~days, one step further each time you get it right and back to the start when you miss it. The number next to \ui{Review} in the menu tells you how many questions are due. A short session every day or two works better than one long session before the exam.

\subsection{Mock exam}
Choose the chapters, the number of questions and the time limit, then start. You get no feedback until you submit; when time runs out, the exam is submitted as it stands. You can flag a question and come back to it. After you submit, you see your score by chapter and question by question. You can also start a mock exam on a single chapter from its \ui{Overview} tab.

\subsection{My report}
The report brings everything together, chapter by chapter: questions practiced, right first time, justification ratings, mock-exam results and predictions on the examples. Readiness is computed from your first answers to the practice questions and from your mock-exam answers only; justification ratings and predictions are shown next to it. Each chapter gets a readiness level:

\begin{center}
\begin{tabularx}{0.86\linewidth}{@{}>{\sffamily\bfseries}l X@{}}
\toprule
\normalfont\textsf{\textbf{Readiness}} & \textsf{\textbf{Meaning}} \\
\midrule
Ready & 80 percent or more of your answers right \\
Review & between 60 and 79 percent right \\
Not yet & below 60 percent right \\
Not enough evidence & you have seen fewer than half of the chapter's questions \\
\bottomrule
\end{tabularx}
\end{center}

The report also lists the questions you missed and have not yet answered correctly. You can save it as a file or print it.

\section{Settings and keeping your progress safe}\label{sg:data}

In \ui{Settings} you can choose the theme (light or dark), a larger text size, the pace of the review (relaxed, standard or intensive) and how many questions a review session holds. You can also set the minimum length of your written justifications and the default length of the mock exam.

Your progress lives in the browser you use, on that computer only. It is lost if you clear the browser's data, and a private window keeps nothing. To keep a copy, or to move to another computer or browser:
\begin{enumerate}
  \item Open \ui{Settings}, then \ui{Your data}, and click \ui{Save progress file}. Keep the file somewhere safe.
  \item On the other computer or browser, open the app, go to the same place and click \ui{Load progress file}.
\end{enumerate}
Save a progress file every week or so, and always before you change computer.

\section{Feedback and error reports}

The app is new, and your reports improve it. Open \ui{Feedback} in the menu on the left. You can tell the instructor how the app works for you, report an error in a question or its solution or in an interactive example, or point out a passage of BHT that you find unclear. Each question card has a \ui{Report} link, and each example a \ui{Report a problem with this example} link, that open the page with the chapter and the item already filled in.

Write what is wrong and what you think it should be, then click \ui{Open in my email program}: the app prepares the message and your email program sends it, so the instructor can reply. If no email program opens, click \ui{Copy the message} and paste it into an email or a Teams chat. Nothing is sent without you.

\section{A routine for each chapter}

The app works best alongside the textbook, not instead of it. For each chapter:
\begin{enumerate}
  \item \textbf{Read the chapter} in BHT first, including its figures and worked examples.
  \item \textbf{Reproduce the examples.} Open each interactive example and, where its introduction says the numbers come from BHT, check that you find them in the text.
  \item \textbf{Predict, then check.} Answer the \ui{Predict first} questions before moving anything. Where you were wrong, go back to the section of BHT the example cites.
  \item \textbf{Explore.} Change one assumption at a time and explain to yourself, in words, why the equilibrium moves the way it does.
  \item \textbf{Answer the questions} in the \ui{Questions} tab, writing real justifications.
  \item \textbf{Do BHT's end-of-chapter exercises}, which the app does not repeat.
  \item \textbf{Review} the due questions every day or two.
  \item \textbf{Take a mock exam} on the chapter, then look at your report. Aim for \emph{Ready} before you move on, or come back to the chapters marked \emph{Review} or \emph{Not yet}.
\end{enumerate}

\section{What each chapter contains}

\begin{center}
\small
\begin{tabularx}{\linewidth}{@{}>{\raggedright\arraybackslash}p{0.30\linewidth} >{\raggedright\arraybackslash}X@{}}
\toprule
\textsf{\textbf{BHT chapter}} & \textsf{\textbf{Interactive examples}} \\
\midrule
1. Why health economics? & The rising health share of GDP; Pooling health risk, and the incentive to hide it \\
\addlinespace
2. Demand for health care & Simple and arc elasticities; Copayments and the effective price; Why a survey misreads demand \\
\addlinespace
3. The Grossman model & An optimal day: the frontier for health and the home good; The MEC curve and the optimal health stock; Aging, education and endogenous death \\
\addlinespace
4. Socioeconomic disparities in health & Disparities in the MEC diagram; One gradient, three causal paths; Selection bias and the Dutch famine \\
\addlinespace
7. Demand for insurance & Expected utility and risk aversion; Fair, unfair, full and partial contracts; Partial-but-fair or full-but-unfair? \\
\addlinespace
8. Akerlof's market for lemons & The market for lemons; The adverse selection death spiral; Quality floors and price ceilings \\
\addlinespace
9. The Rothschild-Stiglitz model & Contracts in the income space; Why no pooling contract survives; The separating equilibrium and when it fails \\
\addlinespace
10. Adverse selection in real markets & The Harvard death spiral; The positive correlation test; Missing early deaths in life insurance data \\
\addlinespace
11. Moral hazard & Price distortion, price sensitivity and social loss; Deductible plus copay; Pauly's trade-off: optimal coverage under moral hazard \\
\addlinespace
15. The health policy conundrum & How full should public insurance be?; Rationing care by price or by queue; Why a country can look inefficient \\
\addlinespace
16. The Beveridge model & Rationing knee replacements by waiting; Global budget or payment by results?; Cost-effectiveness rationing on a fixed budget \\
\addlinespace
17. The Bismarck model & Risk adjustment and cream skimming; Payroll contributions and the private opt-out; Fee schedules and clinical distortions \\
\addlinespace
18. The American model & Wage pass-through and pooling; The tax exclusion for employer premiums; The Medicaid notch \\
\addlinespace
20. Health externalities & Herd immunity and the Pigouvian subsidy; Antibiotic resistance and the Pigouvian tax; The Coase theorem and transaction costs \\
\addlinespace
22. Obesity & Body mass index; Prices, technology and the energy balance; Pooled insurance, moral hazard and a fat tax \\
\bottomrule
\end{tabularx}
\end{center}

Every chapter also has 10 true-or-false and 20 multiple-choice questions: 450 questions in all.

\section{Common questions}

\begin{center}
\small
\begin{tabularx}{\linewidth}{@{}>{\raggedright\arraybackslash}p{0.33\linewidth} >{\raggedright\arraybackslash}X@{}}
\toprule
\textsf{\textbf{Question}} & \textsf{\textbf{Answer}} \\
\midrule
My progress has disappeared. & You probably opened the app in another browser or in a private window, or the browser's data was cleared. Open it in your usual browser; if the data is gone, load your last progress file. \\
\addlinespace
I do not see the correct answer when I am wrong. & The solutions of that chapter are still locked. Your instructor will send the code. \\
\addlinespace
My unlock code does not work. & Check it against the message you received: letters and digits only. If it still fails, tell your instructor which chapter and which version of the file you have. \\
\addlinespace
The \ui{True} and \ui{False} buttons are gray. & Write a longer justification; the counter under the box shows how many words are still needed. \\
\addlinespace
A number in an example is not in BHT. & Read the example's introduction: it says which values come from BHT and which are illustrative. \\
\addlinespace
Chapters 5, 6 and others of BHT are missing. & The app covers only the chapters on the course syllabus. \\
\addlinespace
I found a mistake or something unclear. & Use \ui{Feedback} or the \ui{Report} link next to the question or example. Your feedback shapes the next version. \\
\bottomrule
\end{tabularx}
\end{center}
